% ===========================================================================
%  user/05-projects/00-chapter.tex
% ===========================================================================
\chapter{Projects}\label{cap:projects}

A \voce{project} is a combination of several strategies, with the rule that
decides how much each one weighs and any brakes. It is what you put into
production: running a single strategy is rarely a good idea.

\figuraSchermata{1.0}{figure/schermate/projects.png}
  {The list of saved projects.}
  {projects}

\section{The builder, in four steps}\label{sez:builder}

It opens with \voce{Ensemble builder}, or with \pulsantesec{Use in
Ensemble} from the page of a stream. The four steps at the top always say
where you are.

\subsection{Step 1 --- Strategies}

You choose the streams to combine. You need at least two: with one there is
nothing to combine, and the \pulsante{Next} button stays off until two are
chosen.

\figuraSchermata{1.0}{figure/schermate/composizione.png}
  {Choosing the streams. Next to each one are its Sharpe ratio and return, so
   you do not have to look them up elsewhere.}
  {builder-streams}

\begin{perche}{Do not pick the best three}
The temptation is to take the three highest Sharpe ratios. It is almost
always the mistake: if the three strategies make the same bet, putting them
together reduces nothing --- you are just making it three ways. The
\voce{Corr} column on the home screen is there for this: two strategies with
low correlation, even with modest Sharpe ratios, do better together than
three good strategies that are all alike.
\end{perche}

\subsection{Step 2 --- Weighting}

How to share the capital among the chosen streams. Fourteen methods, grouped
by family.

\figuraSchermata{1.0}{figure/schermate/composizione-pesi.png}
  {The weighting methods, by family. Each one carries its line of
   explanation.}
  {builder-weights}

\begin{center}
\begin{tabular}{@{}L{36mm}L{102mm}@{}}
\toprule
\textbf{Family} & \textbf{The idea} \\
\midrule
\rigaTenue
\voce{Naive} & All the same. \voce{Equal}, an equal share each. Always worth
  trying: it is the benchmark the others are measured against. \\
\voce{Risk / covariance} & Weights by how much each one swings and how much
  they move together. Whatever swings less weighs more. \\
\rigaTenue
\voce{Robust} & As above, but without trusting the estimated correlations
  too much --- on little data they are unstable. \\
\voce{Tail-risk} & Looks at extreme losses instead of the average swing. \\
\rigaTenue
\voce{Return-aware} & Also takes into account how much each one is expected
  to earn, not only how much it risks. \\
\voce{Regime-adaptive} & Changes the weights according to the state of the
  market. \\
\bottomrule
\end{tabular}
\end{center}

\begin{pratica}{Where to start}
\voce{Equal}. Then try \voce{Min-Variance LW} or \voce{ERC} and see whether
the drawdown really falls. If a complicated method does not beat equal
weighting, equal weighting is the answer: it has one parameter fewer to get
wrong.
\end{pratica}

\subsection{Step 3 --- Overlay}

Overlays reduce exposure when something happens. Several can be switched on
together: their effects multiply, and the order in which you choose them
does not change the result.

\begin{center}
\begin{tabular}{@{}L{34mm}L{104mm}@{}}
\toprule
\textbf{Overlay} & \textbf{Reduces exposure when\ldots} \\
\midrule
\rigaTenue
\voce{Vol-Regime} & realised volatility rises. \\
\voce{Jump-Vol} & there is a sudden jump: it reacts within a day. \\
\rigaTenue
\voce{Corr-Gate} & the strategies start moving all together, that is when
  diversification disappears. \\
\voce{Regime-GMM}, \voce{Iso-Stress} & market conditions look like those of
  past periods of stress. \\
\rigaTenue
\voce{Macro-Regime} & the macroeconomic signals worsen. \\
\voce{News-Pressure} & the volume of news goes above normal. \\
\bottomrule
\end{tabular}
\end{center}

\begin{attenzione}{An overlay also takes away gains}
A brake that reduces exposure reduces everything: the losses and the gains.
It is worth looking at the two numbers together --- how much the drawdown
fell, and how much the return fell. Switching on five because \enquote{it
can't hurt} leads to a portfolio that is almost always out of the market.
\end{attenzione}

\subsection{Step 4 --- Costs}

The costs: commissions and slippage. The backtest subtracts them.

\begin{nota}{Raising the costs is the most honest check you can make}
A strategy that stays good with the costs doubled is robust. One that stops
working was earning inside the spread, and on a real account it will not
earn.
\end{nota}

\section{Saving}\label{sez:projects-save}

At the end you give it a name and save it. From then on the project appears
in the list, its performance is what the home screen draws, and --- most
importantly --- it can be chosen on the \voce{Live} page (\vediCap{live}).
